\chapter{Experiment Naming Convention}
\label{appn:B}

All experiment runs follow a consistent naming scheme to ensure reproducibility and traceability of results. Each run is identified by a string of the form:

\begin{center}
\texttt{\{Environment\}\_\{condition\}\_seed\{N\}}
\end{center}

\noindent where \texttt{\{Environment\}} is the Gymnasium environment identifier, \texttt{\{condition\}} describes the agent configuration, and \texttt{\{N\}} is the random seed index. For example:

\begin{center}
\texttt{CartPole-v1\_imagination\_symlog\_seed0}
\end{center}

\noindent The three experimental conditions map to the following labels used in figures and tables throughout this report:

\begin{table}[htbp]
\centering
\caption{Mapping between experiment condition strings and report labels.}
\label{tab:naming_convention}
\begin{tabular}{@{}lll@{}}
\toprule
\textbf{Condition String} & \textbf{Label} & \textbf{Description} \\
\midrule
\texttt{model\_free}            & A & Baseline model-free SAC agent \\
\texttt{imagination}            & B & SAC with world-model imagination \\
\texttt{imagination\_symlog}    & C & SAC with imagination and symlog targets \\
\bottomrule
\end{tabular}
\end{table}